%%%%%%%%%%%%%%%%%%%%%%%%%%%%%%%%%%%%%%%%%
% Masters/Doctoral Thesis 
%
% Template license:
% CC BY-NC-SA 3.0 (http://creativecommons.org/licenses/by-nc-sa/3.0/)
%
%%%%%%%%%%%%%%%%%%%%%%%%%%%%%%%%%%%%%%%%%

% https://www.overleaf.com/latex/templates/template-for-a-masters-slash-doctoral-thesis/mkzrzktcbzfl#.Wj_WWFT1VPW

%----------------------------------------------------------------------------------------
%   Кофигуриране на пакети и опции.
%----------------------------------------------------------------------------------------

\documentclass[12pt,bulgarian,singlespacing,headsepline,oneside,openany]{IICTBASDoctoralThesis}

% Използване на UNICODE.
\usepackage{ucs}
\usepackage[utf8x]{inputenc}

% Рефериране на думи и други обекти.
\usepackage{nameref}

% Рефериране на хипер връзки.
\usepackage{hyperref}

% Кодировка за международни езици.
\usepackage[T1]{fontenc}

% ???
\usepackage{url}

% ???
\usepackage{lipsum}

% Използване на графика.
\usepackage{graphicx}

% Директория в която се намират изображенията.
\graphicspath{{images/}}

% ???
\usepackage{array}

% Автоматично създаване на различните индекси.
\usepackage{imakeidx}

% ???
\usepackage{placeins}

% Автоматично създаване на азбучен указател.
\makeindex[columns=1, title=Азбучен указател, intoc]

%----------------------------------------------------------------------------------------
%   Настройки на страницата.
%----------------------------------------------------------------------------------------

\geometry{
	% Размер на листа.
	paper=a4paper,
	% Вътрешно отместване.
	inner=2.5cm,
	% Външно отместване.
	outer=3.8cm,
	% Отместване за свързване.
	bindingoffset=.5cm,
	% Отместване от горния край.
	top=1.5cm,
	% Отместване от долния край.
	bottom=1.5cm, % Bottom margin
	% Рамка на самата страница.
	%showframe,
}

%----------------------------------------------------------------------------------------
%   Данни за дисертацията.
%----------------------------------------------------------------------------------------

% Заглавие на дисертацията.
\thesistitle{ИЗПОЛЗВАНЕ НА ЕВРИСТИКИ ПРИ ОПТИМИЗАЦИЯ И МАШИННО САМООБУЧЕНИЕ}

% Вид на научната степен.
\degree{"Доктор на науките"} 

% Имена на автора.
\author{\href{http://www.linkedin.com/in/todor-balabanov-2805393}{д-р инж. Тодор Димитров \textsc{Балабанов}}}

% Служебен адрес.
\addresses{ИИКТ-БАН, ул. "акад. Георги Бончев", блок 2, етаж 5, кабинет 514, град София 1113, България} 

% Научна област.
\subject{01.01.12 "Информатика" \\ 4.6 "Информатика и компютърни науки" \\ 4 "Природни науки, математика и информатика"} 

% Ключови думи.
\keywords{} 

% Учебно заведение.
\university{\href{http://bas.bg}{Българска академия на науките}}

% Подразделение.
\faculty{\href{http://iict.bas.bg}{Институт по информационни и комуникационни технологии}} 

% Работна група.
\department{\href{http://iinf.bas.bg}{Моделиране и оптимизация}}

% Параметри на генерирания PDF файл.
\AtBeginDocument{
% Заглавие на PDF файла.
	\hypersetup{pdftitle=\ttitle}
% Автор на PDF файла.
	\hypersetup{pdfauthor=\aname}
% Ключови думи на PDF файла.
	\hypersetup{pdfkeywords=\keywordnames}
}

%----------------------------------------------------------------------------------------
%   Начало на същинския документ.
%----------------------------------------------------------------------------------------

\begin{document}

% Използване на римска номерация за страниците предхождащи същинското изложение.
\frontmatter

% Изчистване на стила за страницата преди да започне същинското изложение.
\pagestyle{plain}

%----------------------------------------------------------------------------------------
%   Заглавна страница.
%----------------------------------------------------------------------------------------

\begin{titlepage}
\begin{center}

\vspace*{.01\textheight}

% Название на учебното заведение.
{\scshape \LARGE 
% Лого на академията.
\includegraphics[width=0.08\linewidth]{logo-bas-bg}
\univname \\[0.5cm] 
% Лого на института.
\includegraphics[width=0.1\linewidth]{logo-iict-bg}
\facname \\[0.5cm] 
\deptname \par}

\vspace{1.5cm}
 
% Автор.
\aname

\vspace{1.5cm}

% Хоризонтална линия.
\HRule \\[0.4cm]

% Заглавие на дисертацията.
{\large \bfseries \ttitle \par}\vspace{0.4cm}

% Хоризонтална линия.
\HRule \\[1.5cm]

% Вид документ.
\textsc{\Large ДИСЕРТАЦИЯ}\\[0.5cm] 
 
\vfill

% Текст изискван от учебното заведение.
\large \textit{за присъждане на научна степен\\ \degreename}\\[1.3cm]

%Област на дисертацията.
\subjectname \vspace{1.5cm}
 
\vfill
 
% Година на предаване.
{\large София \\ 2018}
 
\vfill
\end{center}
\end{titlepage}

%----------------------------------------------------------------------------------------
%   Декларация за достоверност.
%----------------------------------------------------------------------------------------

\begin{declaration}

\vspace{2cm}

Декларирам, че настоящата дисертация съдържа оригинални резултати, получени при проведените от мен научни изследвания. Резултатите, които са получени, описани и/или публикувани от други учени, са надлежно и подробно цитирани в библиографията.

\vspace{0.5cm}

Настоящата дисертация не е представяна в процедури за придобиване на научна степен в друго висше училище, университет или научен институт.

\vspace{2cm}

\hspace{4cm}
Подпис: 

\vspace{0.25cm}

\hspace{5cm}
/ \aname /
\end{declaration}

%----------------------------------------------------------------------------------------
%   Таблицата на съдържанието.
%----------------------------------------------------------------------------------------

\newpage
\tableofcontents

%----------------------------------------------------------------------------------------
%   Списък с фигурите.
%----------------------------------------------------------------------------------------

\newpage
\listoffigures

%----------------------------------------------------------------------------------------
%   Списък с таблиците. 
%----------------------------------------------------------------------------------------

\newpage
\listoftables

%----------------------------------------------------------------------------------------
%   Списък с абревиатурите.
%----------------------------------------------------------------------------------------

\begin{abbreviations}{ll}
\end{abbreviations}

%----------------------------------------------------------------------------------------
%   Основно изложение организирано в глави.
%----------------------------------------------------------------------------------------

% Номерация с арабски цифри за същинската част на дисертацията.
\mainmatter

% Оформление в стил на дисертация.
\pagestyle{thesis}

% По този начин номерацията на подточките е с арабски цифри.
\renewcommand\thesection{\thechapter.\arabic{section}}
\renewcommand\thesubsection{\thesection.\arabic{subsection}}

Time-series Extreme Event Forecasting with Neural Networks at Uber

Предложение за иновативен модел на рекурентна изкуствена неверонна мрежа. Нуждата от подобрения идва от, това че LSTM не дава достатъчно приемливи резултати. Обучението на мрежата става с хетерогенен времеви ред. Целта на изследването е да прогнозира екстремални събития (дни и часове с голямо натоварване). 

Rough Neuron based RBF Neural Networks for Short-Term Load Forecasting

Оптимизация на параметрите в изкуствена невронна мрежа с използването на генетични алгоритми. Изкуствената невронна мрежа се използва за краткосрочна прогноза за енергийно натоварване (електрическа консумация). Представени са „груби“ (rough) неврони в радиално-базисна изкуствена невронна мрежа. Този тип неврони служат за подобряване на обобщаващите функции и за предпазване от преобучение (overfitting). Към всеки неврон във входния и скрития слой се добавят по два входа за „горна“ и „долна“ граница. 

An overview and comparative analysis of Recurrent Neural Networks for Short Term Load Forecasting

Споменават се двупосочно свързани мрежи, където прогнозата не зависи само от миналите периоди, но и от бъдещите периоди. Представен е и алгоритъмът за обучение с обратно разпространение на грешката във времето (Back Propagation Through Time) да „дълбоки“ мрежи. Представен е също метод за обучение на принципа на стохастично градиентно спускане (Stochastic Gradient Descent). Споменават се и мрежи с магистрална архитектура (Highway Network), където информацията преминава през група слоеве без да бъде мониторирана твърде много. Като инструмент за сравнените на прогнозиращи модели първоначално се използва времеви ред на Mackey-Glass. Едно от най-съществените заключения е, че по-ефективна прогноза се постига от по-сложните модели на изкуствени неверонни мрежи. 


%\include{chapters/chapter02} 
%\include{chapters/chapter03}
%\include{chapters/chapter04} 
%\include{chapters/chapter05} 

%----------------------------------------------------------------------------------------
%   Приложения към основното изложени.
%----------------------------------------------------------------------------------------

\appendix

%\include{аppendices/аppendix01}
%\include{аppendices/аppendix02}
%\include{аppendices/аppendix03}

%----------------------------------------------------------------------------------------
%   Списък с използвана литература и източници на информация.
%----------------------------------------------------------------------------------------

\newpage
\begin{thebibliography}{99}
\end{thebibliography}

%----------------------------------------------------------------------------------------
%   Азбучен указател на използваните термини.
%----------------------------------------------------------------------------------------

\newpage
\printindex

\end{document}
